\scanpage{63}
\deleted{algébrique}\inserted{$D$} à coefficients \uline{rationnels}, \deleted{\unsure{dans}}\inserted{sur} $X\times Y$ tel que l'hom. $\deleted{\unsure{a(z)}}a({}^t\gamma(D)):H^{2n-1}(X,\mathbb Q_l)(n-1)\longrightarrow H^1(Y,Q_l)$ soit inverse de $a(\gamma(z))$, ce qui caractérise $D$ comme élément de $\unsure{\operatorname{Cl}(X\times Y)|\operatorname{Cl}X\times Y)|\mathcal C\ell(X\times Y)}\otimes Q$. L'\deleted{hom}élément de $\operatorname{Hom}_{Q}(A_Y,B_X)$ défini par $D$ n'est autre que $u_Z^{-1}$, \deleted{\unsure{à un signe}} Pour qu'on puisse choisir $D$ à coefficients entiers, il faut et il suffit que $u$ soit un isom. \deleted{Pour qu'on} Plus généralement, si $\operatorname{Ker}u$ est de cardinal $N$, \deleted{\unsure{on p}} (Attention, $\operatorname{Ker}u$ en tant que \uline{schéma} en groupes) alors on peut choisir \inserted{$D$} de la forme $D_{\unsure{0}}=\deleted{\unsure{p}}N^{-r}D_0$, $D_0$ étant un div. à coeff. entier, et $r\geq0$, \deleted{\unsure{sss} l'entier $m\geq1$}\inserted{mais pas sous la} forme $\deleted{\unsure{p}}M^{-r}D_0$, \deleted{\unsure{dès|doit}} $\unsure{m|M}$ ne contient pas tous les diviseurs premiers de $n$.

On écrit le diagramme comm
\diagram{\begin{tikzcd}[ampersand replacement=\&,column sep=large,row sep=large]
H^1(Y,Q_l)\arrow[r,"{a(\gamma(z))}"] \& H^{2n-1}(X,Q_l)(n-1)\arrow[d,"\unsure{\wr|2}"] \\
H^1(A_Y,Q_l)\arrow[u,swap,"\wr"]\arrow[r,"u^{(1)}"] \& H^1(B_X,Q_l)
\end{tikzcd}}
où \unsure{toutes} les flèches écrites sont $\mathbb Z$-algébriques, les flèches verticales des isom. \deleted{\unsure{Si}} Pour que
% archive sheet number: 62

